\documentclass[tikz,border=8pt]{standalone}
\usepackage{fontspec}\setmainfont{TeX Gyre Heros}
\usetikzlibrary{arrows.meta}
\definecolor{c64blue}{RGB}{44,41,177}\definecolor{c64lblue}{RGB}{109,106,239}
\definecolor{spr}{RGB}{170,60,60}\definecolor{ref}{RGB}{90,150,90}
\definecolor{gfx}{RGB}{109,106,239}\definecolor{cac}{RGB}{200,150,60}
\definecolor{idle}{RGB}{200,200,200}\definecolor{froz}{RGB}{60,60,60}
\definecolor{wri}{RGB}{200,150,60}
\begin{document}
\begin{tikzpicture}[x=0.24cm,y=0.75cm,font=\footnotesize]
% ======== regle des cycles ========
\foreach \c in {1,5,10,15,20,25,30,35,40,45,50,55,60,63}{
  \node[gray] at (\c-0.5,4.05) {\tiny\c};
}
\draw[gray!60] (0,3.75) -- (63,3.75);
% ======== rangee 1 : acces VIC phase 1 ========
\node[anchor=east] at (-0.6,3.1) {VIC $\phi_1$};
% sprites 3-7 : p-access cycles 1,3,5,7,9 (+s dans les demi-cycles suivants)
\foreach \c/\n in {1/3,3/4,5/5,7/6,9/7}{
  \draw[fill=spr!70,draw=spr!50!black] (\c-1,2.7) rectangle ++(2,0.8);
  \node[white] at (\c,3.1) {\textbf{s\n}};
}
% refresh 11-15
\draw[fill=ref!70,draw=ref!50!black] (10,2.7) rectangle ++(5,0.8);
\node[white] at (12.5,3.1) {\textbf{refresh}};
% g-access 16-55
\draw[fill=gfx!70,draw=gfx!50!black] (15,2.7) rectangle ++(40,0.8);
\node[white] at (35,3.1) {\textbf{g-access (40 octets graphiques)}};
% idle 56-57
\draw[fill=idle,draw=gray] (55,2.7) rectangle ++(2,0.8);
% sprites 0-2 : 58,60,62
\foreach \c/\n in {58/0,60/1,62/2}{
  \draw[fill=spr!70,draw=spr!50!black] (\c-1,2.7) rectangle ++(2,0.8);
  \node[white] at (\c,3.1) {\textbf{s\n}};
}
% ======== rangee 2 : phase 2 (badline) ========
\node[anchor=east] at (-0.6,1.9) {VIC $\phi_2$};
\draw[fill=cac!75,draw=cac!50!black] (14,1.5) rectangle ++(40,0.8);
\node at (34,1.9) {\textbf{c-access (40 pointeurs + couleurs) — Bad Line seulement}};
% ======== rangee 3 : CPU ligne normale ========
\node[anchor=east] at (-0.6,0.7) {6510 (ligne normale)};
\draw[fill=gray!25,draw=gray] (0,0.3) rectangle ++(63,0.8);
\node at (31.5,0.7) {63 cycles disponibles (moins les fenêtres sprites actives)};
% ======== rangee 4 : CPU badline ========
\node[anchor=east] at (-0.6,-0.5) {6510 (Bad Line)};
\draw[fill=gray!25,draw=gray] (0,-0.9) rectangle ++(11,0.8);
\node at (5.5,-0.5) {\tiny 11 cy};
\draw[fill=wri!75,draw=wri!50!black] (11,-0.9) rectangle ++(3,0.8);
\node at (12.5,-0.5) {\tiny écr.};
\draw[fill=froz,draw=black] (14,-0.9) rectangle ++(40,0.8);
\node[white] at (34,-0.5) {\textbf{CPU gelé 40 cycles (AEC bas, 15--54)}};
\draw[fill=gray!25,draw=gray] (54,-0.9) rectangle ++(9,0.8);
\node at (58.5,-0.5) {\tiny 9 cy};
% legende bas
\node[align=left,anchor=north west] at (0,-1.5) {\tiny
Positions des accès : sprites 3--7 = cycles 1--10 ; refresh = 11--15 ; g-access = 16--55 ; sprites 0--2 = 58--63.\quad
Fenêtres sprites : volées seulement si le sprite est \textbf{actif} — sprites 3--7 servent la ligne courante, sprites 0--2 la ligne \textbf{suivante} (BA bas 3 cycles avant ; « écr. » = BA bas, écritures CPU encore possibles).\quad
Sources : [pal.timing] · [Bauer §3.6.3]};
\end{tikzpicture}
\end{document}
